\documentclass[aspectratio=169, 10pt]{beamer}
\usepackage{beamer_theme_cienciadados2026}

\title[Aula 10 - Pandas II: ETL]{Ciência dos Dados I: Análise Exploratória}
\subtitle{Aula 10: Pandas II: Pipeline de Limpeza, Agregação (GroupBy) e Junções (Merge)}
\author{Prof. Dr. Ney Lemke}
\institute{UNESP -- Instituto de Biociências de Botucatu}
\date{Versão 2026}

\begin{document}

\frame{\titlepage}

\begin{frame}{Sumário da Aula}
  \tableofcontents
\end{frame}

\section{Tratamento de Dados Ausentes}

\begin{frame}[fragile]{Estratégias para Missing Data}
  No mundo real, dados faltantes são a regra, não a exceção:
  \begin{itemize}
    \item \textbf{Identificação}: \texttt{df.isna().sum()} reporta o total de nulos por coluna.
    \item \textbf{Descarte}: \texttt{df.dropna(subset=['coluna'])} remove observações incompletas.
    \item \textbf{Imputação}: \texttt{df['coluna'].fillna(valor)} preenche ausências com média, mediana ou constante.
  \end{itemize}
  \begin{lstlisting}
# Imputando pela mediana (robusta a outliers)
mediana = df['pressao'].median()
df['pressao'] = df['pressao'].fillna(mediana)
  \end{lstlisting}
\end{frame}

\section{O Paradigma Split-Apply-Combine}

\begin{frame}[fragile]{Agregações Poderosas com \texttt{groupby}}
  \begin{columns}
    \begin{column}{0.5\textwidth}
      \begin{enumerate}
        \item \textbf{Split}: Divide os dados em grupos com base em critérios.
        \item \textbf{Apply}: Aplica uma função de cálculo a cada grupo isoladamente.
        \item \textbf{Combine}: Reúne os resultados em uma única tabela sumária.
      \end{enumerate}
    \end{column}
    \begin{column}{0.48\textwidth}
      \begin{lstlisting}
# Agregacoes multiplas com .agg()
resumo = df.groupby('setor').agg(
    media_temp=('temperatura', 'mean'),
    desvio_temp=('temperatura', 'std'),
    total_pacientes=('id', 'count')
)
      \end{lstlisting}
    \end{column}
  \end{columns}
\end{frame}

\section{Junção de Tabelas (Joins)}

\begin{frame}[fragile]{Relacionando Bases com \texttt{pd.merge}}
  Semelhante ao comando \texttt{JOIN} do SQL:
  \begin{itemize}
    \item \textbf{Inner Join} (padrão): Retorna apenas registros presentes em ambas as tabelas.
    \item \textbf{Left Join}: Mantém todas as linhas da tabela da esquerda e anexa atributos da direita.
  \end{itemize}
  \begin{lstlisting}
df_final = pd.merge(df_pacientes, df_medicos, on='id_setor', how='left')
  \end{lstlisting}
\end{frame}

\begin{frame}{Resumo da Aula}
  \begin{itemize}
    \item O pré-processamento de dados (ETL) viabiliza modelos e estatísticas confiáveis.
    \item \texttt{groupby} e \texttt{merge} são as operações de maior retorno analítico no Pandas.
    \item \textbf{Prática}: Imputação de dados faltantes, agregações segmentadas e junções de tabelas.
  \end{itemize}
\end{frame}

\end{document}
